% CMPE 220 — Assignment 2: Program Layout & Execution (report)
% Author: Faiq Malik  |  Spring 2026
%
% Figures for \includegraphics{...} live in ../figures/ (compile from reports/).
%
% Build PDF (from repository root, optional output directory):
%   pdflatex -interaction=nonstopmode -output-directory=reports reports/Assignment2_Program_Layout_Report.tex
%   pdflatex -interaction=nonstopmode -output-directory=reports reports/Assignment2_Program_Layout_Report.tex
%
% Edit \githuburl below before submitting.

\documentclass[12pt]{article}

\usepackage[margin=1in]{geometry}
\usepackage[doublespacing]{setspace}
\usepackage{url}

\pagestyle{plain}

\newcommand{\githuburl}{https://github.com/USERNAME/REPONAME}

\begin{document}

% ---------------------------------------------------------------------------
% Title page
% ---------------------------------------------------------------------------
\begin{titlepage}
    \thispagestyle{empty}
    \centering
    \vspace*{1.5in}

    {\LARGE\bfseries CMPE 220 --- System Software}\\[0.35in]
    {\Large Spring 2026}\\[0.75in]

    {\LARGE\bfseries Assignment 2}\\[0.2in]
    {\LARGE\bfseries Program Layout and Execution}\\[1.25in]

    {\large\bfseries Team members}\\[0.2in]
    {\large Faiq Malik}\\[1.5in]

    {\large Date: \today}
\end{titlepage}

\clearpage
\setcounter{page}{1}

% ---------------------------------------------------------------------------
% GitHub repository
% ---------------------------------------------------------------------------
\section*{GitHub Repository}

\noindent
The complete source code, any built artifacts or example binaries included in the repository, demonstration videos, and documentation are maintained at:

\vspace{0.5in}
\begin{center}
\large\url{\githuburl}
\end{center}

\vspace{0.35in}
\noindent
The repository root contains \texttt{cpu16\_core.cpp} (CPU16 assembler and emulator), \texttt{factorial.c} (recursive factorial with \texttt{main} driver), prebuilt images such as \texttt{fact.bin} and \texttt{fib.bin} when provided, demonstration media under \texttt{demos/} (including \texttt{demos/Demo.mp4} and \texttt{demos/Demo\_Recursion\_Memory.mp4}), diagrams under \texttt{figures/} when present, and this report under \texttt{reports/}.

\clearpage

% ---------------------------------------------------------------------------
% Download, compile, and run
% ---------------------------------------------------------------------------
\section*{Download, Compile, and Run}

\subsection*{Download}

\noindent
Clone the repository with Git, or download a ZIP archive from the GitHub page above and extract it to a working directory on macOS or Linux.

\subsection*{Requirements}

\begin{itemize}
  \item C++ compiler with C++17 support (\texttt{g++} recommended).
  \item C compiler with C11 support (\texttt{gcc} recommended) for the factorial program.
  \item Terminal access.
\end{itemize}

\subsection*{Compile}

\noindent
From the repository root directory, build CPU16:

\begin{verbatim}
g++ -std=c++17 -O2 -o cpu16 cpu16_core.cpp
\end{verbatim}

\noindent
This produces the executable \texttt{cpu16}, which implements assembler mode (\texttt{asm}), bare-memory emulator mode (\texttt{emu}), and assemble-and-run mode (\texttt{run}).

\vspace{0.35in}

\noindent
Build the C recursive factorial program:

\begin{verbatim}
gcc -std=c11 -O2 -o factorial factorial.c
\end{verbatim}

\subsection*{Run}

\noindent
Run the C program (enter a non-negative integer when prompted):

\begin{verbatim}
./factorial
\end{verbatim}

\vspace{0.35in}

\noindent
Emulate a factorial program image with memory dump (when \texttt{fact.bin} is present in the project directory, from assembly or distributed with the repository):

\begin{verbatim}
./cpu16 emu fact.bin --base 0x0000 --pc 0x0000 --dump 0x0000 0x01FF
\end{verbatim}

\noindent
This loads the binary into the simulated 64K address space at base \texttt{0x0000}, starts at program counter \texttt{0x0000}, and prints a hex dump of \texttt{0x0000}--\texttt{0x01FF} so code, data, and stack regions are visible while recursive calls use \texttt{CALL}/\texttt{RET}.

\vspace{0.35in}

\noindent
Typical commands when example assembly sources are under \texttt{examples/}:

\begin{verbatim}
./cpu16 run examples/hello.asm
./cpu16 run examples/timer.asm
./cpu16 asm examples/fib.asm -o fib.bin
./cpu16 emu fib.bin --base 0x0000 --pc 0x0100 --dump 0x0100 0x0140
\end{verbatim}

\noindent
The Fibonacci listing uses \texttt{.org 0x0100}; when emulating \texttt{fib.bin}, set the program counter to \texttt{0x0100} as in the command above. If \texttt{fib.bin} lives in the repository root instead of beside the assembly file, use the same arguments with that path, for example:

\begin{verbatim}
./cpu16 emu fib.bin --base 0x0000 --pc 0x0100 --dump 0x0100 0x0140
\end{verbatim}

\noindent
To assemble factorial from source when \texttt{examples/fact.asm} is available, then load and dump memory:

\begin{verbatim}
./cpu16 asm examples/fact.asm -o fact.bin
./cpu16 emu fact.bin --base 0x0000 --pc 0x0000 --dump 0x0000 0x01FF
\end{verbatim}

\subsection*{Video}

\noindent
The repository includes \texttt{demos/Demo.mp4} (software CPU and tool-chain overview) and \texttt{demos/Demo\_Recursion\_Memory.mp4} (function calls, recursion in C and on CPU16, assembly context, and memory layout using emulator dumps).

\clearpage

% ---------------------------------------------------------------------------
% Contributions
% ---------------------------------------------------------------------------
\section*{Team Member Contributions}

\subsection*{Faiq Malik}

\noindent
Contributions include:

\begin{itemize}
  \item \textbf{Recursive C program (\texttt{factorial.c}):} Implemented a recursive factorial function with a base case for $n \leq 1$ and a recursive multiply-and-call step; wrote the \texttt{main} driver that reads user input, validates non-negative input, invokes the recursive function, and prints the result. Added comments tying the program to the stack-layout comparison with the simulated CPU.
  \item \textbf{Software CPU (Assignment~1 carryover, \texttt{cpu16\_core.cpp}):} Maintained the 16-bit CPU implementation: instruction set, two-pass assembler, emulator, memory-mapped UART and timer, and command-line interface (\texttt{run}, \texttt{asm}, \texttt{emu}). Kept code formatted and commented so graders can follow control flow and memory behavior.
  \item \textbf{Program layout demonstration:} Produced or maintained assembled images (for example \texttt{fact.bin}) and documented how to use \texttt{emu} with \texttt{--dump} to inspect memory layout, illustrating how function calls and recursion map to concrete addresses on CPU16.
  \item \textbf{Repository organization:} Structured the repository with a clear layout: sources at the root, \texttt{reports/} for LaTeX PDF sources, \texttt{demos/} for videos, \texttt{figures/} for diagrams, plus binaries and build outputs as needed.
  \item \textbf{Video deliverables:} Recorded \texttt{demos/Demo.mp4} (CPU16 overview) and \texttt{demos/Demo\_Recursion\_Memory.mp4} (function calls, recursion execution, assembly-level discussion, and memory layout).
  \item \textbf{Deliverables:} Repository layout, grader-facing build and run documentation (including this report), demonstration videos as required, supporting assets as required, and this report submitted as PDF for the course.
\end{itemize}

\end{document}
